\section{What the society does on its own}
\label{sec:no-disc}

Nothing in \S\ref{sec:model} refers to an agent's class, so before adding any bias
it is worth asking what this society does unaided. It polarises --- and the
\emph{order} in which it polarises turns out to be controlled by a single ratio.

\subsection{Measuring a society}
\label{sec:order-params}

A society's state is a set of weight vectors and an $N\times N$ matrix of distrust.
To compare states across parameters --- and, in principle, across populations built
from different kinds of agent --- we need statistics that depend on neither the
dimension $K$ nor the parameterisation.

Two pairwise quantities do the work. The \emph{opinion overlap} of two agents is
the cosine of their weight vectors,
$\rho_{IJ}=\hat{w}_I\cdot\hat{w}_J/(\|\hat{w}_I\|\|\hat{w}_J\|)$, which is $+1$
when they agree on every issue and $-1$ when they disagree on every one. The
\emph{trust} a receiver $I$ places in an emitter $J$ is $\eta_{J|I}=1-2\Phi(\hmu)$,
which is $+1$ when $J$ is fully trusted and $-1$ when fully distrusted: the distrust
field mapped onto a bounded, sign-carrying scale. Their correlation,
%
\begin{equation}
  \Rwmu = \bigl\langle (\eta_{I|J}+\eta_{J|I})\,\rho_{IJ} \bigr\rangle,
  \label{eq:rwmu}
\end{equation}
%
averaged over unordered pairs and normalised to $[-1,1]$, asks whether agents trust
those who agree with them.

A correlation cannot distinguish a society split into two coherent camps from one
in which distrust is spread inconsistently, so we also measure social balance in the
sense of \citet{heider1946attitudes} and \citet{cartwright1956structural}. A triple
of agents is balanced in opinion when $b^{I}_{IJK}=\rho_{IJ}\rho_{JK}\rho_{KI}>0$
and balanced in affect when
$b^{A}_{IJK}=\tfrac{1}{2}(\eta_{IJ}\eta_{JK}\eta_{KI}+\eta_{JI}\eta_{IK}\eta_{KJ})>0$;
the aggregates
%
\begin{equation}
  \BI = \bigl\langle b^{I}_{IJK}\bigr\rangle_{(IJK)},
  \qquad
  \BA = \bigl\langle b^{A}_{IJK}\bigr\rangle_{(IJK)}
  \label{eq:balance}
\end{equation}
%
run over all $\binom{N}{3}$ triples. (The subscript $I$ is historical, from the
older name for the opinion sector; it does not index an agent.) A randomly
initialised society has $\BI\approx\BA\approx0$: half its triples are frustrated. A
society that has settled into two internally coherent, mutually opposed blocs has
$\BI=\BA=1$, whatever the blocs are made of --- and a society that cannot settle at
all stays near zero or goes negative. Balance therefore separates \emph{organised}
disagreement from mere disorder, which is the distinction the phase diagram of
\S\ref{sec:results} turns on. Both aggregates are computed in closed form from
matrix products rather than by enumerating triples
(Appendix~\ref{app:closed-form}).

Note that none of \eqref{eq:rwmu} or \eqref{eq:balance} mentions weights, gradients
or $K$. They need only, per pair of agents, a signed alignment of opinion and a
signed trust. Any population in which those two things can be elicited --- whatever
its agents are internally --- can therefore be placed on the same axes, and two
$N\times N$ matrices is a cheap enough measurement to run on a deployed system.
\S\ref{sec:discussion} returns to why a per-agent audit cannot substitute for it.

\subsection{Polarisation, and which kind comes first}

An unbiased society does not converge on consensus. Learning drives each agent away
from the dissonant quadrants of Figure~\ref{fig:modulation}, and the two consonant
quadrants --- agree-and-trust, disagree-and-distrust --- are equally attractive, so
the society organises into two mutually distrustful camps whose membership has
nothing to do with any label. Both balances rise from zero towards one: frustration
is annealed away, and what remains is structured opposition.

Which balance rises first is the interesting part. Figure~\ref{fig:trajectories}
follows the society through the $(\BI,\BA)$ plane from its random initial condition,
at $d=0$, for agendas spanning four orders of magnitude in the ratio
%
\begin{equation}
  \alpha \;=\; \frac{P}{K},
  \label{eq:alpha}
\end{equation}
%
the number of issues in play per dimension of the space they are represented in.
For a simple agenda the path runs up the $\BA$ axis: affective structure forms
almost completely before opinions begin to sort, ending at
$(\BI,\BA)=(0.02,0.97)$ for $\alpha=0.03$. For a complex agenda the path stays
below the diagonal, ending at $(0.97,0.74)$ for $\alpha=333$: opinions separate
first and distrust follows. The crossover is at $\alpha\approx1.7$, where the
trajectory lies along the diagonal.

\begin{figure}[t]
  \centering
  \includegraphics[width=0.46\linewidth]{figures/agenda_trajectories}
  \caption{\textbf{Which polarisation comes first is set by the agenda.}
  Trajectories in the $(\BI,\BA)$ plane with no discrimination ($d=0$), one curve
  per agenda complexity $\alpha=P/K$, markers equally spaced in time --- they crowd
  at the end because the annealed dynamics slows, as \S\ref{sec:variables} requires.
  Simple agendas run above the diagonal (distrust forms first), complex ones below
  it (disagreement first), crossing over at $\alpha\approx1.7$.}
  \label{fig:trajectories}
\end{figure}

So the question of whether groups come to dislike each other because they disagree,
or come to disagree because they dislike each other, has a parameter attached to it
rather than an answer: narrow agendas --- few issues, functioning as symbols --- put
affect first, broad ones put opinion first, and the two are regimes of one system
rather than rival accounts. The same ratio explains, in \S\ref{sec:results}, why one
of the two discriminatory phases exists only for narrow agendas.
